\section{Related Work}
\label{s:related_work}

%\subsection{Traditional methods in residual disease prediction}
%To predict the residual disease of ovarian cancer, researchers manually analyze or design models based on the data generated from various preoperative examinations.
Researchers have put a lot of effort into predicting residual disease in ovarian cancer\cite{chi2006optimal, ferrandina2009role, shim2015nomogram}, including manually analyzing or designing models based on data generated from various preoperative examinations.
In the early stages, doctors systematically observe the whole abdomen through laparoscopy to predict the residual disease\cite{testa2012ultrasound, van2019laparoscopy}. However, Chi et al.\cite{chi2004ten} have highlighted the potential risks of invasive examination and associated complications.

%After that, numerous researchers have focused on non-invasive methods.
%Several studies perform multivariate analysis with preoperative clinical data which includes the physical, hematological, genetic, and other indicators of the patient.
%Laios et al.\cite{laios2020predicting, laios2020247} combined age, Body Mass Index (BMI), the timing of surgery, and other variables to construct a k-Nearest Neighbor (k-NN) model for prediction.
%Berchuck et al.\cite{berchuck2004prediction} analyzed the association of residual disease with genes, and screened 120 genes to develop a predictive model.

In recent years, a growing number of researchers have shifted their focus to non-invasive methods for predicting residual disease in ovarian cancer. Many studies have employed multivariate analysis using preoperative clinical data, including physical, hematological, genetic, and other patient indicators. 
For instance, Laios et al.\cite{laios2020predicting, laios2020247} used age, Body Mass Index (BMI), surgery timing, and other variables to develop a k-Nearest Neighbor (k-NN) model for prediction. Meanwhile, Berchuck et al.\cite{berchuck2004prediction} examined the association of residual disease with genes and screened 120 genes to construct a predictive model.
%Through stepwise logistic regression analysis, Shah et al.\cite{shah2018combining} proposed that serum microRNA and serum CA-125 levels could be employed as predictive markers.
%Shah et al.\cite{shah2018combining} proposed using serum CA-125 levels and serum microRNA levels as predictive markers through stepwise logistic regression analysis.
Another approach proposed by Shah et al.\cite{shah2018combining} is the use of serum CA-125 levels and serum microRNA levels as predictive markers through stepwise logistic regression analysis. 
Moreover, in recent years, preoperative medical image analysis with deep learning methods~\cite{shi2021darnet} has become an increasingly influential method for predicting residual disease. Medical image data, such as CT and Magnetic Resonance Imaging (MRI), are visual and informative.
For instance, Barber et al.\cite{barber2021natural} utilized preoperative CT reports to develop a Natural Language Processing (NLP) model.
Li et al.\cite{li2021noninvasive} employed Radiomics\cite{gillies2016radiomics} features extracted from MRI with clinical data. 
Last but not least, Peritoneal Carcinomatosis Index (PCI), an imaging indicator based on disease size and location within the peritoneal cavity, has been widely applied to predicting residual disease \cite{llueca2018prediction, engbersen2019mri, jonsdottir2021peritoneal, li2022diffusion}.
%\subsection{Deep-learning methods of ovarian cancer}

%Although no studies have been performed to predict the residual disease of ovarian cancer, deep-learning methods have been performed in other tasks related to ovarian cancer, such as diagnosis\cite{saida2022diagnosing}, histotype classification\cite{farahani2022deep} and recurrence prediction\cite{wang2019deep}. 

Although no studies have been conducted specifically to predict residual disease in ovarian cancer using deep-learning methods, such methods have been applied to related tasks, such as diagnosis\cite{saida2022diagnosing}, histotype classification\cite{farahani2022deep}, and recurrence prediction\cite{wang2019deep}. 
%These studies match or outperform the performance of experienced doctors, showing the potential of deep-learning methods and medical image analysis.
Notably, these studies have demonstrated that deep-learning methods and medical image analysis can match or even outperform experienced doctors, highlighting their potential for clinical use.

% todo
